% Copyright (c) 2020 by Daniel R. Driver.
% !TEX encoding = UTF-8 Unicode
% !TEX TS-program = XeLaTeX

\documentclass[titlepage]{article}

% This document presumes a file structure and set of inputs that are
% available at: git@github.com:danieldriver/syllabi.git

\newcommand\policy{../policy}
\newcommand\incl{../includes}
\ProvidesFile{variables.tex}[2022/09/11 v2.3 -- Syllabus variables]

\usepackage{xspace} % make manual spaces (like \mycmd\ ) unnecessary
\usepackage{xifthen} % provides \isempty test

% variables for internal use
\newcommand\prof{}
\newcommand\pdegree{}
\newcommand\pphone{}
\newcommand\pemail{}
\newcommand\poffice{}
\newcommand\phours{}
%
\newcommand\ccode{}
\newcommand\ctitle{}
\newcommand\cseries{}
\newcommand\cversion{}
\newcommand\csemester{}
\newcommand\cmeetson{}
\newcommand\cmeetsat{}
\newcommand\cmeetsin{}
\newcommand\cwebsite{}
\newcommand\cdescrip{}
\newcommand\cprereqs{}
\newcommand\edobject{}

% in case of fully online, BTh, or undergraduate courses - https://tex.stackexchange.com/a/5896
\newif\ifonline
\newcommand\Int[2]{\ifonline#1\else#2\fi}
\newif\ifbth
\newcommand\bth[2]{\ifbth#1\else#2\fi}
\newif\ifUG
\newcommand\UG[2]{\ifUG#1\else#2\fi}

% commands for setting variables in the preamble
\newcommand\professor[2][PhD]{
  \renewcommand\pdegree{#1\xspace}
  \renewcommand\prof{#2\xspace}}
\newcommand\phone[1]{
  \renewcommand\pphone{\addfontfeatures{Numbers=Monospaced}#1\xspace}}
\newcommand\email[1]{
  \renewcommand\pemail{\href{mailto:#1}{#1}\xspace}}
\newcommand\officehours[2][AST Library/Teams]{
  \renewcommand\poffice{#1\xspace}
  \renewcommand\phours{\href{https://outlook.office.com/bookwithme/user/b97b147a2e5c43b4a0e200c2cb630b77@astheology.ns.ca?anonymous}{#2}\xspace}%
  }
%
\newcommand\coursecode[2][1.0]{
  \renewcommand\cversion{#1\xspace}
  \renewcommand\ccode{#2\xspace}}
\newcommand\coursetitle[2][]{
  \ifthenelse{\isempty{#1}}%
    {}% do nothing if #1 is empty, else:
    {\renewcommand\cseries{#1\\[1ex]}}
  \renewcommand\ctitle{#2\xspace}}
\newcommand\semester[1]{
  \renewcommand\csemester{#1\xspace}}
\newcommand\meets[3]{
  \newcommand\AM{\textsc{am}}
  \newcommand\PM{\textsc{pm}}
  \renewcommand\cmeetson{#1\xspace}
  \renewcommand\cmeetsat{\Int{From 9:00 \AM}{#2}\xspace}
  \renewcommand\cmeetsin{\Int{\href{https://smu.brightspace.com/d2l/login}{Brightspace}}{#3}\xspace}}
\newcommand\website[2][]{
  \renewcommand\cwebsite{
  \ifthenelse{\isempty{#1}}%
    {\href{https://#2}}% copy url if #1 is empty, else
    {\href{#1}}% use specified url and print #2
  {#2}\xspace}}
\newcommand\cdescription[2][BF 1001; RM 1000 or GTRS 6000]{
  \renewcommand\cprereqs{#1}
  \renewcommand\cdescrip{#2\par}}
\newcommand\objectives[1]{
  \renewcommand\edobject{#1\par}}


\onlinetrue % \Int{true}{false}
\coursecode[3.5]{BF 1001}
\coursetitle[Biblical Foundations]{Hebrew Bible/Old Testament}

% Taught as RLGS 1013 in:
%   - Fall 2008
%   - Winter 2009
%   - Fall 2009
%   - Winter 2010
%   - Fall 2010
%   - Winter 2011 (Fall 2011 was covered by Joel Lohr in a course release)
%   - Winter 2012
%   - Fall 2012
%   - Winter 2013
% Taught as BSTH 1013 in:
%   - Fall 2013
%   - Winter 2014
%   - Fall 2014
%   - Winter 2015
%   - Fall 2015
%   - Winter 2016
% Taught as BF 1001 in:
%   - Fall 2016
%   - Fall 2017
%   - Fall 2017(Int - online section run parallel to campus course)
%   - Fall 2018
%   - Winter 2020(Int)

\professor{Daniel R. Driver}
\phone{902-222-0930}
\email{ddriver@astheology.ns.ca}
\officehours{Thursdays, 12:30--2:00 \PM}

\semester{Winter Term 2020}
\meets{\Int{Fridays}{Mondays}}% \meets{on}{at}{in}
      {1:00--3:30 \PM}% if online, it prints "From 9:00 AM"
      {Classroom 3}% if online, it prints "Brightspace"
\website{danieldriver.com}
\cdescription[none. The course is required for HB courses at or above the 3000-level. It is open to Continuing Education participants]{% copy from the current Academic Calendar; [] for prereqs
	The First Testament of Christian Scripture, also called the Hebrew
	Bible or Old Testament, is a fundamental part of Christian tradition
	and durable rule of faith and practice. Students in this course will
	therefore be introduced to historical and literary data important
	for understanding the origins of the Hebrew Bible and its subsequent
	function as Scripture in a variety of cultural and religious
	contexts.

	The broad sweep of biblical tradition will be presented through a
	survey of representative books from the Law, the Prophets, and the
	Writings. To help foster an ecumenical outlook, primary readings
	from the Bible will be supplemented by case studies of classic
	Jewish and Christian readings of biblical texts. Students will
	develop their exegetical skills by studying these examples, and so
	learn to appreciate the diverse literary, canonical, cultural,
	historical, hermeneutical, and theological elements involved in
	biblical interpretation. In addition, the course will consider some
	ways that the Hebrew Bible/Old Testament can feature in the practice
	of ministry and in the spiritual practice of faith communities.
}% end of course description
\objectives{% recall Bloom's taxonomy: http://www.celt.iastate.edu/teaching/RevisedBlooms1.html
	By the end of the course students should be able to:
		name major Old Testament people and events;
		locate a few important biblical sites on a map;
		give key dates for Israel's history and summarize the succession of superpowers in the ancient Near Eastern political theatre from Egypt to Greece;
		recognize and cite examples of most genres of biblical literature;
%		defend a decision to pronounce or circumlocute the divine name;
		understand the general shape of the Masoretic Text tradition and differentiate it from other canonical orders;
		report on parallel and divergent material across the Law and the Prophets, such as the uses of the Divine Name or the rationales for sabbath observance;
		classify and begin to evaluate a variety of ancient and modern approaches to the Bible.

	Students should also be able to identify settings in which the
	Scriptures of Israel are read (notably the synagogue, church, and
	academy), employ terminology appropriate to these communities,
	recognize where their own biographies place them in relation to the
	Hebrew Bible/Old Testament and its uses, and monitor and test their
	individual attitudes and assumptions. They should be able to extend
	their awareness of the Bible's contemporary readers to the Bible's
	long history of reception. Finally, students should begin to infer
	what Jesus meant in speaking of ``the law of Moses, the prophets,
	and the psalms'' (Luke 24:44), and so learn to hear claims about New
	Testament fulfilment of scripture in light of the unique voice that
	the Old Testament retains along side of the New in Christian
	Scripture.
}% end of learning objectives

\ProvidesFile{preamble.tex}[2014/09/02 v1.1 -- Syllabus preamble]

% basic typography
\usepackage{fontspec,realscripts}
\defaultfontfeatures{Ligatures=TeX}
\setmainfont[%
  Numbers={Proportional,OldStyle},
  SmallCapsFeatures={RawFeature=+c2sc;+smcp,LetterSpace=1}]{Scala OT}
\newfontfamily\lining[Numbers={Lining}]{Scala OT}
\setsansfont[%
  SmallCapsFeatures={RawFeature=+c2sc;+smcp,LetterSpace=1}]{Scala Sans Pro}
\setmonofont[Scale=MatchLowercase]{Menlo}
\usepackage{sectsty}
\allsectionsfont{\sffamily\scshape}
\frenchspacing
\setlength{\emergencystretch}{3em} % prevent overfull lines

% custom font size and leading
\renewcommand\tiny{\fontsize{6}{9}\selectfont}
\renewcommand\scriptsize{\fontsize{7}{10}\selectfont}
\renewcommand\footnotesize{\fontsize{8}{10}\selectfont}
\renewcommand\small{\fontsize{9}{11}\selectfont}
\renewcommand\normalsize{\fontsize{10}{12}\selectfont}% base size
\renewcommand\large{\fontsize{11}{13}\selectfont}
\renewcommand\Large{\fontsize{12}{14}\selectfont}
\renewcommand\LARGE{\fontsize{16}{18}\selectfont}% "course syllabus \\ semester" benefits from more lead
\renewcommand\huge{\fontsize{20}{22}\selectfont}
\renewcommand\Huge{\fontsize{24}{26}\selectfont}

% layout packages: page, logo, tables
\usepackage[scale={0.6,0.8},
            xetex]{geometry}
\usepackage{graphicx}
\usepackage{array}     % allow insertions of column styling with >{}
\usepackage{booktabs}  % elegant horizontal rules in tables

% custom macros for a session count in the schedule of readings
\newcounter{session}
\newcounter{columns}
\newcounter{spancolumns}
\newcounter{midcolumns}
\newcounter{courseunit}
\newcommand\setcolumncount[2][0]{ % optionally set count to other than 0,
  \setcounter{session}{#1}        % e.g. to -1, or to a standing count
  \setcounter{columns}{#2}
  \setcounter{spancolumns}{#2-1}
  \setcounter{midcolumns}{#2-2}}
\newcommand\sessioncount{\stepcounter{session}\arabic{session}}
\newcommand\sessionskip[1]{\multicolumn{1}{@{}r@{ }}{#1}}
\newcommand\unit[1]{\multicolumn{\thecolumns}{c}{%
  \scshape\stepcounter{courseunit}\roman{courseunit}. #1}}
\newcommand\noclass[1]{\sessionskip{---.}&\multicolumn{\themidcolumns}{@{}l}{\itshape No Class: #1}}
\newcommand\rarr{\Large\char"EA31\hspace*{0.5pt}} % pointer finger w/ reminder tie in Scala OT
\newcommand\reminder[2]{\sessionskip{\rarr}&\multicolumn{\thespancolumns}{@{}l}{#1 \hfill #2}}

% color to match Tyndale's branding
\usepackage[usenames]{xcolor}
% predefined: black, white, red, green, blue, cyan, magenta, yellow
\definecolor{TyndaleURLs}{HTML}{0062A0} % links on tyndale.ca
\definecolor{TyndaleBlue}{cmyk}{1,1,0,.32}
\definecolor{TyndaleGold}{cmyk}{0,.27,1,0}
\definecolor{TyndaleRed}{cmyk}{0,1,.99,.04}
\definecolor{TyndaleBlack}{cmyk}{0,0,0,1}
\definecolor{TyndaleGreen}{cmyk}{.45,0,1,.24}
\definecolor{TyndaleOrange}{cmyk}{0,.79,1,0}
\definecolor{TyndaleAqua}{cmyk}{.47,0,.24,0}
\definecolor{TyndaleYellow}{cmyk}{.03,.03,.35,0}

% metadata (assumes a host of definitions are made in the main file)
\usepackage[setpagesize=false,     % leave this to geometry
            hyperfootnotes=false,  % fragile and distracting
            xetex]{hyperref}
\hypersetup{breaklinks=true,       % allow link text to break across lines
            colorlinks=true,       % colorlinks resets pdfborder to 0 0 0
            urlcolor=TyndaleURLs,  % for external links
            linkcolor=TyndaleRed,  % for normal internal links
            citecolor=TyndaleGold, % for bibliographical citations in text
            pdfauthor={\prof},
            pdftitle={\ccode: \ctitle},
            pdfsubject={Tyndale UC, \csemester},
            pdfcreator={github.com/danieldriver/syllabus}}
\urlstyle{same}                    % don't use monospace font for urls

% custom footlines
\usepackage{fancyhdr}
\pagestyle{fancy} % turn it on
\fancyhf{}        % reset everything
\renewcommand{\headrulewidth}{0pt} % remove header line as well
\lfoot{\sffamily\scshape\footnotesize\ctitle, \csemester}
\rfoot{\sffamily\scshape\footnotesize\prof\quad\thepage}

% language support: polyglossia needs to go last
\usepackage{polyglossia}
\setdefaultlanguage{english}
\setotherlanguage{hebrew}
\newfontfamily\hebrewfont[Script=Hebrew]{SBL BibLit}

% protect positioning of margin table in policy/grades -- needs to come after polyglossia
\usepackage{marginfix}

% gratuitous with custom title page, but useful as a fallback
\title{\ccode: \ctitle}
\author{\professor}
\date{\semester}


\begin{document}
\ProvidesFile{title.tex}[2013/09/06 v1.0 -- Syllabus title page]

\begin{titlepage}
  \begin{center}

    \begin{minipage}{\textwidth}
      \parbox[t]{0.5\textwidth}{
        \mbox{}\\[-\baselineskip] % dummy line to align parboxes
        \includegraphics[width=0.5\textwidth]{.syllabus/includes/TyndaleUC}}
      \hfill
      \parbox[t]{0.4\textwidth}{
        \raggedleft\LARGE\sffamily Course Syllabus\\
        \semester}
    \end{minipage}

    \vfill

    {\sffamily\textsc{\LARGE\MakeLowercase\coursecode}\\[2ex]
      \bfseries\itshape\Huge\coursetitle}

    \vfill

    \begin{tabular}{>{\bfseries}rl>{\bfseries}rl}
      \toprule
      Instructor & \professor                   & Course  & \liningnums{\coursecode}~\texttt{(v\version)} \\
      \midrule
      Phone      & \texttt{\phone}              & Meets   & \meeton                                       \\
      Email      & \href{mailto:\email}{\email} & Time    & \meetat                                       \\
      Office     & \officenumber                & Room    & \meetin                                       \\
      Hours      & \officehours                 & Website & \href{http://\website}{\website}              \\
      \bottomrule
    \end{tabular}

    \vfill

    \begin{description}
      \item[Commuter Hotline]
        Class cancellations due to inclement weather or illness will
        be announced on the commuter hotline at \texttt{416.226.6620
        x2187}. Alternately, weather cancellation information is posted
        at \href{http://tyndale.ca/weather}{tyndale.ca/weather}.
      \item[MyTyndale.ca]
        This course may have materials stored on its website, such as
        handouts or readings that may be needed in order to complete
        assignments. Students are responsible for checking these course
        pages on a regular basis. Here, too, students are able to view
        their grades throughout the semester. For more information see
        Section~\ref{mytyndale}, below.
      \item[Mail]
        Students are responsible for information communicated through
        their campus mailboxes and student e-mail accounts. A mailbox
        directory hangs beside the mailboxes. For more information
        contact the Registrar's office.
    \end{description}

  \end{center}

  \section{Course Description}
  \label{description}

  \emph{From the Academic Calendar:} \cdescription

\end{titlepage}
\setcounter{page}{2} % count the title page as page 1


  \section{Learning Objectives}
  \label{objectives}

  \edobject

\section{Required Texts \& Materials}
\label{texts}

The following texts are required. Students are strongly encouraged to
purchase their own copies. Library copies that are not reference works
have been placed on a 2-hour reserve.

\begingroup
\renewcommand{\section}[2]{}% temporarily remove the section heading
\begin{thebibliography}{Irenaeus}% use the longest item in the bibliography

	\bibitem[NRSV]{nrsv} M.\,D. Coogan, ed.
    \emph{New Oxford Annotated Bible with Apocrypha: NRSV}. 5th ed.
    Oxford: Oxford University Press, 2018.
    ISBN 978-0190276072.

	\bibitem[HBFB]{hbfb} J. Kaminsky and J. Lohr.
	\emph{The Hebrew Bible for Beginners: A Jewish and Christian Introduction}.
	Nashville: Abingdon Press, 2015.
	ISBN 978-1426775635.

	\bibitem[Heschel]{heschel} Abraham Heschel.
	\emph{The Sabbath: Its Meaning for Modern Man}.
	New York: Farrar, Straus and Giroux, 1951 (repr. 2005).
	ISBN 978-0374529758.

	\bibitem[Irenaeus]{irenaeus} St Irenaeus of Lyons.
	\emph{On the Apostolic Preaching}.
	Trans. John Behr.
	Crestwood, NY: St Vladimir’s Seminary Press, 1997.
	ISBN 978-0881411744.

\end{thebibliography}
\endgroup

An good alternative study Bible is the NJPS: Adele Berlin and Marc Zvi
Brettler, eds., \emph{The Jewish Study Bible: Second Edition} (Oxford:
Oxford University Press, 2014). \href{https://danieldriver.com/courses/bf-1001/}{See
my website for links to required and recommended books, and a discussion
of options for a Bible.}

\section{Supplementary Texts}
\label{supplementary}

Supplementary readings will be recommended throughout the semester and
either placed on reserve or made available through the course website.
See \autoref{bib}. Give this material good effort and attention. It
represents some of the best contemporary readers of the Bible, and you
will not want to miss their hermeneutical and theological insights.

If you are looking to build your biblical studies library, I recommend
Konrad Schmid's \emph{A Historical Theology of the Hebrew Bible}
\cite{ks19} for a recent introduction to critical scholarship that goes
beyond \cite{hbfb}. Brevard Childs \cite{bc79, bc92} remain useful, too.
Also, the following reference works are worth owning and consulting.
\cite{rlgs} in particular contains sound advice on core skills like
reading religious texts, writing essays and reviews, revising essays,
making oral presentations, and learning languages.

\begingroup
\renewcommand{\section}[2]{}% temporarily remove the section heading
\begin{thebibliography}{Making Sense}% use the longest item in the bibliography

	\bibitem[Making Sense]{rlgs} Northey, Margot, Bradford A. Anderson, and Joel N. Lohr.
	\emph{Making Sense in Religious Studies: A Student's Guide to Research and Writing}.
	3rd ed. Toronto: Oxford University Press, 2019. ISBN 978-0199026838.

	\bibitem[SBL2]{sbl2} Collins, Billie Jean, et al.
	\emph{The SBL Handbook of Style}.
	2nd ed. Atlanta: SBL Press, 2014. ISBN 978-1589839649.
	Designed to augment \href{http://www.chicagomanualofstyle.org/home.html}{\emph{Chicago Style}}
	(the standard at AST), there is also a free
	\href{https://www.sbl-site.org/assets/pdfs/pubs/SBLHSsupp2015-02.pdf}{Student Supplement for SBL2}.

\end{thebibliography}
\endgroup


\section{Course Outline}
\label{outline}

We will adhere to the schedule in \autoref{schedule} as closely as
possible, though the professor reserves the right to adjust it to suit
the needs of the class.

\newcommand\HBFB[1]{\cite[#1]{hbfb}}

\setcolumncount{5}% set up \sessioncount, \unit{}, \noclass{}, and \reminder{memo}{date} macros
\begin{table}[htbp]% set to `p' to put the schedule on its own page
  \centering
  \begin{tabular}{>{\sessioncount.}r@{ }lllr}% make sure the column config agrees with \setcolumncount
	\toprule
	\sessionskip{\textbf{\S}.}&\textbf{Primary (chapters)}&\textbf{Secondary (pgs)}&\textbf{Supplement (author)}&\textbf{Due}\\
	\midrule

	\unit{Law / Torah / Pentateuch} \\

		& Genesis 1--11        & \HBFB{1--26}    & Davis \cite{ed03}, Seitz  \cite{cs96} & \Int{10}{} Jan. \\
		& Genesis 12--50       & \HBFB{27--64}   & Levenson                  \cite{jl12} & \Int{17}{} Jan. \\
		& Exodus 1--34         & \HBFB{65--75}   & McGinnis                  \cite{cm12} & \Int{24}{} Jan. \\
		& Leviticus 1--16      & \HBFB{77--83}   & Anderson                  \cite{ga17} & \Int{31}{} Jan. \\
		& Deuteronomy          & \HBFB{85--99}   & \cite{heschel} or Moberly \cite{wm13} & \Int{ 7}{} Feb. \\
	\reminder{A first short paper is \textbf{due} by the end of the sixth week of class}{14 Feb.}    \\ [1ex]

	\unit{Prophets: Former \textit{\&} Latter} \\

		& Joshua, Judges       & \HBFB{103--121} & Trible                    \cite{pt84} & \Int{14}{} Feb. \\
	\noclass{AST Reading Week}                                                           & \Int{21}{} Feb. \\
		& 1 \& 2 Samuel        & \HBFB{123--134} & Chapman                   \cite{sc16} & \Int{28}{} Feb. \\
		& 1 \& 2 Kings         & \HBFB{135--143} & Alter                     \cite{ra81} & \Int{ 6}{} Mar. \\
%	\noclass{AST Closed on Monday}                                                       & \Int{2 }{} Mar. \\
		& Isaiah 1--12, 36--40 & \HBFB{145--168} & \cite{irenaeus} or Davis  \cite{ed14} & \Int{13}{} Mar. \\
		& Hosea--Jonah         & \HBFB{169--184} & Nogalski          \cite{jn07a, jn07b} & \Int{20}{} Mar. \\
	\reminder{A second short paper is \textbf{due} by the end of the eleventh week of class}{27 Mar.}    \\ [1ex]

	\unit{Writings: Wisdom \textit{\&} Praise} \\

		& Psalms 1--8, 41, 51, 89–90 & \HBFB{187--202} & deClaissé-Walford  \cite{ndw14} & \Int{27}{} Mar. \\
		& Ecclesiastes         & \HBFB{203--246} & Dell \cite{kd13} or Fox   \cite{mf04} & \Int{ 3}{} Apr. \\ [1ex]

	\reminder{End of Term: Final marks are due for all courses}{14 Apr.} \\

	\bottomrule
  \end{tabular}
  \caption{Schedule of Readings}
  \label{schedule}
\end{table}

See the AST website for a list of other \href{http://www.astheology.ns.ca/students/academic-dates.html}{important dates}.

\section{Evaluation}
\label{evaluation}

The grade structure for \ccode has the following elements.

\begin{enumerate}

	\item \textbf{Reading prompts} will be given throughout the
	semester. They are low-risk writing exercises designed to encourage
	careful reading of the assigned material, and to help explore its
	significance. Individual posts are not marked, though points may be
	deducted from your final grade if the exercise is not taken
	seriously. In general, prompt responses submitted online should not
	exceed 300 words. Try to post the in the morning of class days,
	after you do the reading but before you listen to the lecture. You
	are invited (but not strictly required) to respond to and interact
	with classmates on the forums through the rest of the day.

	\item Two \textbf{short papers} will facilitate student engagement
	with the art of biblical interpretation. One is keyed to the
	exposition of a model, the other to a short biblical text. Each
	should be about 2,500 words in length. Either paper (your choice) is
	due at the end of the sixth week of class; the other one is due at
	the end of the eleventh week.

	\begin{enumerate}

		\item A \textbf{review essay} invites student reflection on a
		model work of biblical interpretation. The short books by
		\cite{heschel} and \cite{irenaeus} are the primary options,
		though with permission you may also select one of the
		supplementary articles. Note that a review is not the same thing
		as a report. Devote the first half of the paper to a summary the
		interpretation or argument under review, and be mindful that
		traditional interpretation requires special attention to
		context. (Irenaeus' reading of Isaiah or the Jewish tradition
		around sabbath observance are not the same thing as an academic
		article on the Book of the Twelve.) Devote the second half of
		the paper to critical analysis and evaluation of your chosen
		example. Be fair, but do not fail to take a position. The paper
		needs to develop a \textbf{thesis}. See me and \cite[Chs 3, 5–7,
		11]{rlgs} for guidance.

		\item An \textbf{exegetical essay} provides an opportunity for
		direct work with the biblical text. The first task is to
		identify an appropriate text. Select a suitably short passage
		from the HB/OT. Then, conduct an analysis and explication of it.
		Interact with \cite{hbfb} and at least two other commentators.
		Advance a \textbf{thesis} that relates to the text itself. See
		me and \cite[Chs 3, 5, 8, 11]{rlgs} for guidance.

	\end{enumerate}

\Int{
	\item Each student will \textbf{make an online presentation} on one
	of the twelve supplementary articles. Be creative! It may take the
	form of a podcast or online video, and should be sharable by URL
	(try SoundCloud, YouTube, or Vimeo). The presentation should last
	15–20 minutes. It must begin with a brief (5~min) summary of the
	article. The balance of the time should be spent helping other
	readers elaborate, reflect on, test, challenge, or extend the main
	ideas presented there.
}{
	\item Each student will lead a \textbf{seminar discussion} of one
	chapter from the HB/OT, freely chosen from a week's primary reading.
	The discussion should last about 25~minutes. It needs to begin with
	a brief (5~min) \textbf{outline} of the chapter. Outlines must be
	original, not borrowed from commentators. They should be circulated
	as handouts and explained orally. The balance of the time should be
	spent helping the class discover, elaborate, reflect on, test, or
	otherwise explore the meaning of the passage. Critical to success
	here is the ability to ask good questions of the Bible.
}

\end{enumerate}

The breakdown for the semester's total work is shown in
\autoref{grade-dist}.

\begin{table}[htbp]
  \centering
  {\lining
  \begin{tabular}{lr}
    \toprule
    Reading Prompts  & 25\% \\
    Review Essay     & 25\% \\
    Exegetical Essay & 25\% \\
    \Int{Online Presentation}{Seminar Discussion} & 25\% \\
    \bottomrule
  \end{tabular}}
  \caption{Distribution of Grades}
  \label{grade-dist}
\end{table}

%AST's \href{http://www.astheology.ns.ca/students/resources.html}{Academic
%Calendar} provides guidelines and detailed criteria for academic
%assessment. Marks are assigned by letter grade using these benchmarks.

\ProvidesFile{grades.tex}[2023/09/06 v2.5 -- Course policy]

AST's \href{https://www.astheology.ns.ca/students/resources.html}{Academic
Calendar} provides guidelines and criteria for academic assessment. Marks are
assigned by letter grade using these benchmarks. Note that graduate and
undergraduate programs are currently assessed on a different scales at AST (see
\autoref{grade-syst}). For graduate students, a grade of C is deemed
``Unsatisfactory,'' whereas for undergraduate (BTh) students, a C is
``Satisfactory'' and the grade of D is possible for ``Marginal'' performance.
Other grades have different definitions, too.

\begin{table}[htbp]
 \centering
 {\lining\footnotesize
 \begin{tabular}{lll|ll}
   \toprule
               & \multicolumn{2}{c|}{Graduate Courses}   & \multicolumn{2}{c}{Undergraduate Courses} \\
               \cline{2-5}
   A+          & 94--100    & Exceptional               & 90--100    & Excellent                 \\
   A           & 87--93     & Outstanding               & 85--89     & Excellent                 \\
   A\char"2212 & 80--86     & Excellent                 & 80--84     & Excellent                 \\ [1ex]
   B+          & 77--79     & Good                      & 75--79     & Good                      \\
   B           & 73--76     & Acceptable                & 70--74     & Good                      \\
   B\char"2212 & 70--72     & Marginally Acceptable     &            &                           \\ [1ex]
   C           & 60--69     & Unsatisfactory            & 60--69     & Satisfactory              \\
   D           &            &                           & 50--59     & Marginal                  \\
   F           & 0--59      & Failure                   & 0--49      & Failure                   \\
   FP          & 0          & Failure due to Plagiarism & 0          & Failure due to Plagiarism \\
   \bottomrule
 \end{tabular}}
 \caption{Outline of Grades at AST}
 \label{grade-syst}
\end{table}

\ProvidesFile{other.tex}[2018/0/03 v2.4 -- Course policy]

\section{Other Course Policy}
\label{policy}

Late work will not be accepted, except in genuinely extenuating
circumstances. Students must submit something before the deadline if
they wish to receive credit. Unless I state otherwise, written
assignments are to be uploaded by 11:59 \PM\ on the date indicated.

Essay submissions must be typewritten, double-spaced, and formatted as
PDFs. They should be free from error. In this course they should follow
either \emph{The SBL Handbook of Style} (refer to the online
\href{https://www.sbl-site.org/assets/pdfs/pubs/SBLHSsupp2015-02.pdf}{
Student Supplement}), or \emph{The Chicago Manual of Style}, per AST's
policy for major papers. As a reminder, AST also upholds an Inclusive
Language Policy.

\href{http://www.eerdmans.com/Pages/Item/59043/Commentary-Statement.aspx}{Plagiarism},
if \href{https://www.theguardian.com/world/2013/feb/09/german-education-minister-quits-phd-plagiarism}{detected},
will result in failure of the course.

Students should request permission to record a class or lecture. If
permission is granted, or if recordings are provided (as in the case of
an online course), I stipulate that all recordings be for personal use
only. They may not be shared or distributed.

If you have abilities or disabilities that require modifications to the
assessment process or other aspects of this course, please advise the
course instructor as soon as possible.

Finally, I encourage the conscientious use of laptops, tablets, and
other technology in my classes. In classroom settings, realize that, as
\href{http://dx.doi.org/10.1016/j.compedu.2012.10.003}{cognitive
psychologists have demonstrated}, ``laptop multitasking hinders
classroom learning for both users and nearby peers.'' Do your part to
foster an environment of open dialogue by honouring the presence of your
classmates. In online settings, consider both the physical environment
in which you choose to work and the virtual environment that you help
create through your participation in various forums. Let your engagement
in this course be marked by rigour and charity alike.

\section{Bibliography}
\label{bib}

Literature on the Bible is vast. The works listed here have been
selected for clarity, insight, and theological alertness. Some will be
distributed as supplementary texts, as per \autoref{supplementary}. All
of them model compelling, highly competent ways of reading the HB/OT.

\begingroup
\renewcommand{\section}[2]{}% temporarily remove the section heading
\begin{thebibliography}{99}

\bibitem{ra81} Alter, Robert. “The Techniques of Repetition.” Pages 88–113 in \emph{The Art of Biblical Narrative}. New York: Basic Books, 1981.

\bibitem{ga01} Anderson, Gary A. “Biblical Origins and the Fall.” Pages 197–210 in \emph{The Genesis of Perfection: Adam and Eve in Jewish and Christian Imagination}. Louisville: Westminster John Knox, 2001.

\bibitem{ga17} Anderson, Gary A. “Apophatic Theology: The Transcendence of God and the Story of Nadab and Abihu.” Pages 3–22 in \emph{Christian Doctrine and the Old Testament: Theology in the Service of Biblical Exegesis}. Grand Rapids: Baker Academic, 2017.

\bibitem{bz09} Ben Zvi, Ehud and James D. Nogalski. \emph{Two Sides of a Coin: Juxtaposing Views on Interpreting the Book of the Twelve / the Twelve Prophetic Books}. Analecta Gorgiana 201. Piscataway, NJ: Gorgias, 2009.

\bibitem{sc16} Chapman, Stephen B. “1 Samuel 1–12.” Pages 71–119 in \emph{1 Samuel as Christian Scripture: A Theological Commentary}. Grand Rapids: Eerdmans, 2016.

\bibitem{bc79} Childs, Brevard S. \emph{Introduction to the Old Testament as Scripture}. Philadelphia: Fortress, 1979.

\bibitem{bc92} Childs, Brevard S. \emph{Biblical Theology of the Old and New Testaments: Theological Reflection on the Christian Bible}. Minneapolis: Fortress, 1992.

\bibitem{bc04} Childs, Brevard S. \emph{The Struggle to Understand Isaiah as Christian Scripture}. Grand Rapids: Eerdmans, 2004.

\bibitem{ed03} Davis, Ellen F. “Teaching the Bible Confessionally in the Church.” Pages 9–26 in \emph{The Art of Reading Scripture}. Edited by Ellen F. Davis and Richard B. Hays. Grand Rapids: Eerdmans, 2003.

\bibitem{ed14} Davis, Ellen F. \emph{Biblical Prophecy: Perspectives for Christian Theology, Discipleship, and Ministry}. Louisville: Westminster John Knox, 2014.

\bibitem{ndw14} deClaissé-Walford, Nancy L. “The Meta-Narrative of the Psalter.” Pages 363–76 in \emph{The Oxford Handbook of the Psalms}. Edited by William P. Brown. Oxford: Oxford University Press, 2014.

\bibitem{kd13} Dell, Katherine J. “Ecclesiastes as Wisdom: Consulting Early Interpreters.” Pages 9–36 in \emph{Interpreting Ecclesiastes: Readers Old and New}. Winona Lake, IN: Eisenbrauns, 2013.

\bibitem{mf04} Fox, Michael V. \emph{Ecclesiastes: The Traditional Hebrew Text with the New JPS Translation}. The JPS Bible Commentary. Philadelphia: Jewish Publication Society, 2004.

\bibitem{ah51} Heschel, Abraham Joshua. \emph{The Sabbath: Its Meaning for Modern Man}. New York: Farrar, Straus and Giroux, 1951.

\bibitem{il97} Irenaeus of Lyons. \emph{On the Apostolic Preaching}. Translated by John Behr. Popular Patristics Series 17. St Vladimir’s Seminary Press, 1997.

\bibitem{bj13} Janowski, Bernd. \emph{Arguing with God: A Theological Anthropology of the Psalms}. Louisville: Westminster John Knox, 2013.

\bibitem{jl85} Levenson, Jon D. \emph{Sinai and Zion: An Entry into the Jewish Bible}. Minneapolis: Winston, 1985.

\bibitem{jl12} Levenson, Jon D. “The Test.” Pages 66–112 in \emph{Inheriting Abraham: The Legacy of the Patriarch in Judaism, Christianity, and Islam}. Princeton: Princeton University Press, 2012.

\bibitem{nm09} MacDonald, Nathan. “Israel and the Old Testament Story in Irenaeus’s Presentation of the Rule of Faith.” \emph{Journal of Theological Interpretation} 3.2 (2009): 281–98.

\bibitem{cm12} McGinnis, Claire Mathews. “The Hardening of Pharaoh’s Heart in Christian and Jewish Interpretation.” \emph{Journal of Theological Interpretation} 6.1 (2012): 43–64.

\bibitem{wm92} Moberly, R. W. L. \emph{The Old Testament of the Old Testament: Patriarchal Narratives and Mosaic Yahwism}. Minneapolis: Fortress, 1992.

\bibitem{wm13} Moberly, R. W. L. “A Love Supreme.” Pages 7–40 in \emph{Old Testament Theology: Reading the Hebrew Bible as Christian Scripture}. Grand Rapids: Baker Academic, 2013.

\bibitem{jn07a} Nogalski, James D. “Reading the Book of the Twelve Theologically.” \emph{Interpretation} 61.2 (2007): 115–22.

\bibitem{jn07b} Nogalski, James D. “Recurring Themes in the Book of the Twelve: Creating Points of Contact for a Theological Reading.” \emph{Interpretation} 61.2 (2007): 125–36.

\bibitem{db89} Patrick, Dale. “Studying Biblical Law as a Humanities.” \emph{Semeia} 45 (1989): 27–47.

\bibitem{ks19} Schmid, Konrad. \emph{A Historical Theology of the Hebrew Bible}. Translated by Peter Altmann. Grand Rapids: Eerdmans, 2019. ISBN 978-0802876935.

\bibitem{cs96} Seitz, Christopher R. “Old Testament or Hebrew Bible? Some Theological Considerations.” \emph{Pro Ecclesia} 5.3 (1996): 292–303.

\bibitem{cs18} Seitz, Christopher R. “‘Can We Read This Book?’ Reader Response-ability.” Pages 51–68 in \emph{The Elder Testament: Canon, Theology, Trinity}. Waco, TX: Baylor University Press, 2018.

\bibitem{gs92} Sheppard, Gerald T. “Theology and the Book of Psalms.” \emph{Interpretation} 46.2 (1992): 143–55.

\bibitem{pt84} Trible, Phyllis. “An Unnamed Woman: The Extravagance of Violence.” Pages 65–91 in \emph{Texts of Terror: Literary-Feminist Readings of Biblical Narratives}. Overtures to Biblical Theology. Philadelphia: Fortress, 1984.

\end{thebibliography}
\endgroup

\end{document}
